\pdfoutput=1
\documentclass[11pt]{article}

\usepackage[a4paper,margin=1.05in]{geometry}
\usepackage{amsmath,amssymb,amsthm,mathtools}
\usepackage{enumitem}
\usepackage[colorlinks=true,linkcolor=blue,citecolor=blue,urlcolor=blue]{hyperref}

\numberwithin{equation}{section}

\newtheorem{theorem}{Theorem}[section]
\newtheorem{lemma}[theorem]{Lemma}
\newtheorem{proposition}[theorem]{Proposition}
\newtheorem{corollary}[theorem]{Corollary}
\newtheorem{conjecture}[theorem]{Conjecture}
\theoremstyle{definition}
\newtheorem{definition}[theorem]{Definition}
\newtheorem{remark}[theorem]{Remark}

\newcommand{\R}{\mathbb{R}}
\newcommand{\St}{\operatorname{St}}
\newcommand{\Gr}{\operatorname{Gr}}
\newcommand{\tr}{\operatorname{tr}}
\newcommand{\diag}{\operatorname{diag}}
\newcommand{\rank}{\operatorname{rank}}
\newcommand{\conv}{\operatorname{conv}}
\newcommand{\spec}{\operatorname{spec}}
\newcommand{\smin}{\sigma_{\min}}
\newcommand{\lmin}{\lambda_{\min}}
\newcommand{\cE}{\mathcal E}
\newcommand{\cA}{\mathcal A}

\title{Sharp Extremals and Active-Set Evidence for the Goreinov--Tyrtyshnikov--Zamarashkin Hypothesis in the First Open Case}
\author{Pavel Osinenko}
\date{\today}

\begin{document}
\maketitle

\begin{abstract}
We study the first open case $(n,k)=(6,3)$ of the
Goreinov--Tyrtyshnikov--Zamarashkin row-selection hypothesis.  In projector
coordinates the problem asks whether every rank-three orthogonal projector
$P\in\R^{6\times6}$ has a $3\times3$ principal submatrix with least eigenvalue
at least $1/6$.  We do not prove the hypothesis.  Instead we give a
computer-assisted exact analysis of the known equality configurations.  First,
the bound is tight: there are explicit projectors with
\[
  \max_{|T|=3}\lambda_{\min}(P_{TT})=\frac16.
\]
Second, six weighted series--parallel graphic-matroid extremals and one further
out-of-family extremal are sharp nonsmooth local minima of the max-eigenvalue
functional on $\Gr(3,6)$.  The sharpness certificates are exact rational or
number-field computations after positive diagonal changes of coordinates.
They imply that all seven configurations are isolated points of the equality
set.  We also record a symmetry-reduced active-set census and numerical KKT
screens over the active-set orbits; no sub-threshold critical point was found.
Finally, for one ten-active equality stratum we give a Groebner/CAD cascade
which eliminates its positive-dimensional residual determinant branches and
classifies the remaining finite residue in the chart.  We also include exact
finite-branch computations for the two size-twelve equality strata and a
zero-dimensional chart computation for the thirteen-active out-of-family
stratum.  For the most persistent structured size-six low-active survivor, we
also give an exact finite-branch exclusion of the full four-parameter ansatz.
As background, we include a self-contained Perron-free proof of the weighted
Frobenius inequality behind the known $k=2$ case.
We also record the known volume-sampling relaxation
\[
  \max_{|I|=k}\sigma_{\min}(A_I)^2\ge \frac{1}{k(n-k)+1},
\]
which gives the concrete bound $F(P)\ge 1/10$ in the case $(6,3)$.
For the first open case we prove a stronger relaxed theorem:
\[
  F(P)\ge \frac19,
\]
i.e. the GTZ bound with relaxation factor $\alpha=3/2$.
The remaining global problem is isolated precisely as the exclusion of critical
points with value below $1/6$.
\end{abstract}

\tableofcontents

\section{Introduction}

Let
\[
  \St(n,k)=\{A\in\R^{n\times k}: A^TA=I_k\}.
\]
The Goreinov--Tyrtyshnikov--Zamarashkin hypothesis predicts that for every
$A\in\St(n,k)$ one can choose $k$ rows forming a $k\times k$ submatrix $A_I$ with
\[
  \smin(A_I)\ge \frac1{\sqrt n}.
\]
The cases $k=1$ and $k=2$ are known, and by duality the case $k=n-2$ follows from
the case $2$.  Thus $(n,k)=(6,3)$ is the first self-dual open case.

For $A\in\St(6,3)$ set
\[
  P=AA^T,\qquad F(P)=\max_{|T|=3}\lambda_{\min}(P_{TT}).
\]
The right action $A\mapsto AQ$, $Q\in O(3)$, does not change $P$, so the natural
domain is the Grassmannian
\[
  \Gr(3,6)=\{P=P^T=P^2,\ \tr P=3\}.
\]
The case $(6,3)$ is the assertion
\begin{equation}\label{eq:gtz63}
  F(P)\ge \frac16\qquad\text{for all }P\in\Gr(3,6).
\end{equation}

This paper is deliberately a partial-result paper.  It does not prove
\eqref{eq:gtz63}.  Its purpose is to record exact local information at the
boundary of the conjectured inequality and to make the remaining global problem
as concrete as possible.

\subsection*{Main contributions}

\begin{enumerate}[leftmargin=*]
\item We give a self-contained Perron-free proof of a weighted Frobenius
inequality for $k=2$, implying the usual good-pair step in the core case.
\item We record a standard volume-sampling corollary which gives the relaxed
bound $F(P)\ge 1/10$ for all $P\in\Gr(3,6)$.
\item We prove the stronger relaxed bound $F(P)\ge1/9$ in the case $(6,3)$.
The proof combines a leverage deletion using the known $(5,3)$ case with a
shifted determinant-sum identity on the high-leverage core.
\item We rederive the tight equality projectors in the weighted graphic-matroid
family and verify exactly that $F=1/6$ there.
\item We give exact sharpness certificates at all six projectors in this
series--parallel census.  In particular the equality projectors are not flat:
$F$ grows linearly in every nonzero Grassmann tangent direction.
\item We record a seventh equality projector, with leverage pattern
$(5/14)^{\times3},(9/14)^{\times3}$, outside the weighted series--parallel census,
and certify its sharpness exactly.
\item We prove the general implication ``sharp equality point implies isolated''
for the equality set
\[
  \cE=\{P\in\Gr(3,6):F(P)=1/6\}.
\]
Thus the seven certified projectors are isolated equality points.
\item We enumerate active-set orbits under row relabelling and run numerical KKT
screens over these orbits.  These screens found no sub-threshold critical point.
\item For the hardest observed size-six low-active survivor, we identify a
four-parameter structured ansatz.  Exact characteristic-zero Groebner and FGLM
computations reduce its numerically relevant branch to finite lexicographic
systems, and exact real-algebraic sign checks exclude all real points on that
branch from satisfying the active and inactive threshold inequalities.
\item For the active set of the base extremal $P(S(e,e,e),e,e,e)$, we record an
exact CAD-style cascade: component relations plus active positive
semidefiniteness reduce the positive-dimensional determinant locus to a finite
infeasible endpoint, while the remaining finite open residue is exactly the
known 32-point sign family in the chart.
\item For the thirteen-active out-of-family active set, the determinant chart
ideal is already zero-dimensional of degree $32$, and its lex basis gives
exactly the corresponding 32-point sign family.
\item For the two remaining size-twelve active-set orbits in the known equality
census, the exact determinant ideals have only the two possible real branch
values $u=\pm1/9$ for the inverse Gram determinant.  Since $u=1/d(Z)>0$ on the
real chart, the positive branch is forced; on that branch both ideals are
zero-dimensional of degree $32$ and give exactly their known sign families.
\end{enumerate}

\subsection*{Status conventions}

Statements called \emph{proved} below are either conventional arguments or exact
finite computations with rational/algebraic arithmetic, implemented in the
accompanying scripts listed in Section~\ref{sec:repro}.  Floating-point searches
are labelled numerical evidence and are not used as proofs.

\section{Projector formulation and elementary reductions}

For $P=AA^T$ write $\ell_i=P_{ii}$ and $c_{ij}=P_{ij}$.  If $T\subset[6]$ has
cardinality $3$, then
\[
  P_{TT}=A_TA_T^T,
\]
so $\lambda_{\min}(P_{TT})=\smin(A_T)^2$.  Hence the original row-selection
bound is exactly \eqref{eq:gtz63}.

\begin{proposition}[Volume-sampling relaxed bound]\label{prop:volume-relaxed}
Let $1\le k<n$ and $A\in\St(n,k)$.  Then there is a $k$-row set
$I\subset[n]$ such that
\[
  \sigma_{\min}(A_I)^2\ge \frac{1}{k(n-k)+1}.
\]
Equivalently, for $P=AA^T$,
\[
  \max_{|I|=k}\lambda_{\min}(P_{II})\ge \frac{1}{k(n-k)+1}.
\]
In particular, for $(n,k)=(6,3)$ one has the unconditional relaxed bound
$F(P)\ge 1/10$.
\end{proposition}

\begin{proof}
Set $X=A^T\in\R^{k\times n}$, so $XX^T=I_k$.  Choose a size-$k$ column set
$S$ by volume sampling, i.e. with probability proportional to
$\det(X_SX_S^T)$.  The volume-sampling expectation formula of
Derezi\'nski and Warmuth~\cite[Theorem 4]{DerezinskiWarmuth} gives
\[
  \mathbb E\bigl[(X_SX_S^T)^{-1}\bigr]\preceq (n-k+1)(XX^T)^{-1}
  =(n-k+1)I_k,
\]
with equality when the volume-sampling distribution has full support.  Taking
traces, some sampled set $S$ satisfies
\[
  \tr\bigl((X_SX_S^T)^{-1}\bigr)\le k(n-k+1).
\]
Let $0<\mu_1\le\cdots\le\mu_k\le1$ be the eigenvalues of $X_SX_S^T$.
The upper bound $\mu_j\le1$ follows from $X_SX_S^T\preceq XX^T=I_k$.
Therefore $\mu_j^{-1}\ge1$ for $j\ge2$, and
\[
  \frac{1}{\sigma_{\min}(X_S)^2}
  =\mu_1^{-1}
  \le k(n-k+1)-(k-1)
  =k(n-k)+1.
\]
Since $X_S=A_S^T$, the singular values are the same as those of $A_S$.
\end{proof}

\begin{lemma}[Case-A reduction]\label{lem:caseA}
Let $1\le k<n$, and let $A\in\St(n,k)$.  If some row has leverage
$\ell_i\le 1/n$, then the validity of $\mathrm{GTZ}(n-1,k)$ implies the desired
GTZ conclusion for this particular matrix $A$.
\end{lemma}

\begin{proof}
Rotate columns so that row $i$ is $(b,0,\ldots,0)$ with $b^2=\ell_i$.  Delete
that row and multiply the first column by $(1-\ell_i)^{-1/2}$.  The resulting
matrix $\widetilde B$ lies in $\St(n-1,k)$.  If $J$ is a $k$-row subset of
$\widetilde B$ and $\varphi(J)$ is the corresponding subset of the original
matrix, then
\[
  A_{\varphi(J)}=\widetilde B_J
  \diag\bigl(\sqrt{1-\ell_i},1,\ldots,1\bigr)Q^T
\]
for an orthogonal $Q$.  Therefore
\[
  \smin(A_{\varphi(J)})
  \ge \sqrt{1-\ell_i}\,\smin(\widetilde B_J).
\]
If $\smin(\widetilde B_J)\ge1/\sqrt{n-1}$, then
\[
  \smin(A_{\varphi(J)})
  \ge \sqrt{1-\frac1n}\,\frac1{\sqrt{n-1}}
  =\frac1{\sqrt n}.
\]
\end{proof}

\begin{remark}
Lemma~\ref{lem:caseA} does not imply that proving $(6,3)$ proves all $(n,3)$.
It only removes rows of leverage at most $1/n$.  To prove all $(n,3)$ by
induction, one must still prove the core case, where every leverage is $>1/n$,
at every intermediate size.
\end{remark}

\begin{theorem}[A relaxed \texorpdfstring{\(\alpha=3/2\)}{alpha=3/2} bound]\label{thm:one-ninth}
For every $P\in\Gr(3,6)$,
\[
  F(P)=\max_{|T|=3}\lambda_{\min}(P_{TT})\ge \frac19.
\]
\end{theorem}

\begin{proof}
Write $P=AA^T$ with $A\in\St(6,3)$, and let $\ell_i=P_{ii}$.

First suppose some leverage satisfies $\ell_i\le4/9$.  Repeat the deletion and
column-rescaling argument from Lemma~\ref{lem:caseA}.  The resulting matrix
lies in $\St(5,3)$.  The case $(5,3)$ is known by duality from the $k=2$
theorem, so there are three rows of the rescaled matrix with squared least
singular value at least $1/5$.  Scaling back to the original matrix gives
\[
  \sigma_{\min}(A_J)^2\ge \frac{1-\ell_i}{5}\ge \frac19.
\]

It remains to treat the high-leverage core, where $\ell_i>4/9$ for all $i$.
Then every triple $T$ has
\[
  \tr(P_{TT})=\sum_{i\in T}\ell_i>\frac43.
\]
Let the eigenvalues of $P_{TT}$ be
$0\le\lambda_1\le\lambda_2\le\lambda_3\le1$.  If
$\lambda_2\le1/9$, then
\[
  \tr(P_{TT})\le 1+\frac19+\frac19=\frac{11}{9}<\frac43,
\]
a contradiction.  Hence $\lambda_2>1/9$ for every triple.  With $t=1/9$, this
means
\[
  \lambda_{\min}(P_{TT})\le t
  \quad\Longleftrightarrow\quad
  \det(P_{TT}-tI_3)\le0
\]
throughout the high-leverage core.

Assume for contradiction that every triple is bad at level $t=1/9$.  Then
\[
  \det(P_{TT}-tI_3)\le0
  \qquad\text{for all } |T|=3.
\]
On the other hand, summing the characteristic-polynomial coefficients over all
triples gives a positive number.  Indeed, if $e_r(P_{TT})$ denotes the $r$-th
elementary symmetric polynomial of the eigenvalues of $P_{TT}$, then
\[
  \det(P_{TT}-tI_3)=e_3(P_{TT})-t e_2(P_{TT})+t^2e_1(P_{TT})-t^3.
\]
Since $P$ is a rank-three projector,
\[
  \sum_{|T|=3}e_3(P_{TT})=1,\qquad
  \sum_{|T|=3}e_2(P_{TT})=4\binom32=12,\qquad
  \sum_{|T|=3}e_1(P_{TT})=\binom52\tr P=30.
\]
Therefore
\[
  \sum_{|T|=3}\det(P_{TT}-tI_3)
  =1-12t+30t^2-20t^3.
\]
At $t=1/9$ this is
\[
  1-\frac{12}{9}+\frac{30}{81}-\frac{20}{729}
  =\frac7{729}>0,
\]
contradicting the nonpositivity of every summand.  Thus at least one triple has
$\lambda_{\min}(P_{TT})>1/9$ in the high-leverage core, and the theorem follows.
\end{proof}

\section{A Perron-free Frobenius proof for the case \texorpdfstring{\(k=2\)}{k=2}}

The case \(k=2\) is known by the theorem of Sengupta--Pautov.  We record here
an independent proof of the weighted pair inequality that gives the good-pair
step in the core case.  The proof is included because it is short, uses only
Frobenius identities and one Rayleigh bound, and clarifies which part of the
\(k=2\) argument does not presently extend to triples.

Let \(A=[a\ b]\in\St(n,2)\), and write the \(i\)-th row as
\[
  r_i=(a_i,b_i),\qquad
  \ell_i=\|r_i\|^2,\qquad
  c_{ij}=\langle r_i,r_j\rangle,\qquad
  p_{ij}=a_ib_j-b_ia_j.
\]
Thus \(\sum_i\ell_i=2\), \(\sum_{i<j}p_{ij}^2=1\), and
\[
  \ell_i\ell_j=c_{ij}^2+p_{ij}^2.
\]
Define the pair bracket
\[
  b_{ij}
  =p_{ij}^2-\frac{\ell_i+\ell_j}{n}+\frac1{n^2}
  =\left(\ell_i-\frac1n\right)\left(\ell_j-\frac1n\right)-c_{ij}^2.
\]
Equivalently, \(b_{ij}\) is the characteristic polynomial of the pair block
\(P_{\{i,j\}\{i,j\}}\) evaluated at \(1/n\).

\begin{theorem}[Weighted Frobenius inequality]\label{thm:k2-frobenius}
For every \(A\in\St(n,2)\),
\[
  \sum_{i<j}p_{ij}^2 b_{ij}\ge0.
\]
\end{theorem}

\begin{proof}
Set \(z_i=a_i+\mathrm i b_i\) and \(\zeta_i=z_i^2\).  Since the two columns of
\(A\) are orthonormal,
\[
  \sum_i\zeta_i=0,\qquad \sum_i|\zeta_i|=2.
\]
Let \(w_i=(\operatorname{Re}\zeta_i,\operatorname{Im}\zeta_i)\), let \(W\) be
the \(n\times2\) matrix with rows \(w_i\), and put
\[
  K=WW^T,\qquad L=\ell\ell^T,\qquad Q=\sum_i\ell_i^2.
\]
Then \(K\mathbf 1=0\), and
\[
  K_{ij}=c_{ij}^2-p_{ij}^2,\qquad
  L_{ij}=c_{ij}^2+p_{ij}^2.
\]
Hence \(p_{ij}^2=\frac12(L-K)_{ij}\) off the diagonal, and the diagonal entries
of \(L-K\) vanish.  Therefore
\[
  \sum_{i<j}p_{ij}^4=\frac18\|L-K\|_F^2.
\]
Using \(\sum_{j\ne i}p_{ij}^2=\ell_i\), one also has
\[
  \sum_{i<j}p_{ij}^2 b_{ij}
  =\sum_{i<j}p_{ij}^4-\frac Qn+\frac1{n^2}
  =\frac18\|L-K\|_F^2-\frac Qn+\frac1{n^2}.
\]

It remains to show \(\|L-K\|_F^2\ge 8Q/n-8/n^2\).  Let
\(\mu_1\ge\mu_2\ge0\) be the two eigenvalues of \(W^TW\).  Then
\(\mu_1+\mu_2=Q\) and
\[
  \|K\|_F^2=\mu_1^2+\mu_2^2,\qquad
  \langle L,K\rangle=\ell^TK\ell.
\]
With \(v=\ell-\frac2n\mathbf1\), the identity \(K\mathbf1=0\) gives
\(\ell^TK\ell=v^TKv\), while \(\|v\|^2=Q-4/n\).  The Rayleigh bound for the
positive semidefinite matrix \(K\) gives
\[
  v^TKv\le\mu_1\|v\|^2=\mu_1\left(Q-\frac4n\right).
\]
Consequently
\[
\begin{aligned}
  \|L-K\|_F^2-\frac{8Q}{n}+\frac8{n^2}
  &=
  Q^2-2\ell^TK\ell+\mu_1^2+\mu_2^2-\frac{8Q}{n}+\frac8{n^2}  \\
  &\ge
  Q^2-2\mu_1\left(Q-\frac4n\right)+\mu_1^2+\mu_2^2
  -\frac{8Q}{n}+\frac8{n^2}  \\
  &=2\left(\mu_2-\frac2n\right)^2\ge0,
\end{aligned}
\]
where the last equality uses \(\mu_1=Q-\mu_2\).  This proves the claim.
\end{proof}

\begin{corollary}[Good pair in the core case]\label{cor:k2-good-pair}
For every \(A\in\St(n,2)\) there is a nonsingular pair \(i<j\) with
\[
  b_{ij}\ge0.
\]
If in addition \(A\) is in the core region \(\ell_i>1/n\) for all \(i\), then
this pair satisfies
\[
  \lambda_{\min}\bigl(P_{\{i,j\}\{i,j\}}\bigr)\ge\frac1n.
\]
\end{corollary}

\begin{proof}
The weights \(p_{ij}^2\) are nonnegative and sum to \(1\).  If every nonsingular
pair had \(b_{ij}<0\), then the weighted sum in
Theorem~\ref{thm:k2-frobenius} would be negative, a contradiction.

For the final assertion, let \(\lambda_1\le\lambda_2\) be the two eigenvalues
of \(P_{\{i,j\}\{i,j\}}\).  Since \(b_{ij}\) is its characteristic polynomial at
\(1/n\), the inequality \(b_{ij}\ge0\) says that \(1/n\) is not strictly between
\(\lambda_1\) and \(\lambda_2\).  In the core region
\(\lambda_1+\lambda_2=\ell_i+\ell_j>2/n\).  Therefore \(1/n\) cannot lie above
both eigenvalues, and hence \(\lambda_1\ge1/n\).
\end{proof}

\begin{remark}
Theorem~\ref{thm:k2-frobenius} is a pair theorem.  It proves the key core-case
step for \(k=2\), and the usual Case-A induction supplies the boundary rows.
It does not by itself produce a good triple when \(k=3\).  In fact, the natural
volume-weighted analogues of this averaged inequality fail in higher rank, so
the \((6,3)\) problem requires the separate active-set and semialgebraic
analysis developed below.
\end{remark}

\section{The nonsmooth active-set problem}

The functional
\[
  F(P)=\max_{|T|=3}g_T(P),\qquad g_T(P)=\lambda_{\min}(P_{TT}),
\]
is a maximum of twenty eigenvalue functions.  At a point where all active least
eigenvalues are simple, $g_T$ is differentiable for active $T$ and
\[
  dg_T(P)[\dot P]=v_T^T\dot P_{TT}v_T,
\]
where $v_T$ is a unit least-eigenvector of $P_{TT}$.

\begin{definition}
The active set at $P$ is
\[
  \cA(P)=\{T: g_T(P)=F(P)\}.
\]
If the active least eigenvalues are simple, we call $P$ Clarke-stationary for
$F$ if
\[
  0\in\conv\{dg_T(P):T\in\cA(P)\}
\]
after restriction to the tangent space $T_P\Gr(3,6)$.
\end{definition}

\begin{proposition}[KKT reduction]\label{prop:kkt}
The inequality \eqref{eq:gtz63} holds if and only if no Clarke-stationary point
of $F$ on $\Gr(3,6)$ has value below $1/6$.
\end{proposition}

\begin{proof}
If a sub-threshold point exists, $F$ attains a sub-threshold global minimum on
the compact manifold $\Gr(3,6)$, and such a minimizer is Clarke-stationary.  The
converse is immediate.
\end{proof}

Thus the whole remaining global problem can be phrased as:
\[
  \text{exclude all critical points of }F\text{ with }F<1/6.
\]

\section{The weighted graphic-matroid equality projectors}

We recall the construction used for the six series--parallel equality
projectors.  Let $G$ be a connected graph with six edges, choose a reduced
incidence matrix $B_{\rm red}$, and let $W=\diag(w_1,\ldots,w_6)$ with
$w_i>0$.  Define
\[
  \widetilde M=W^{1/2}B_{\rm red}^T,\qquad
  P=\widetilde M(\widetilde M^T\widetilde M)^{-1}\widetilde M^T.
\]
Then $P$ is a rank-three orthogonal projector when $G$ has four vertices.  The
leverage of an edge is
\[
  \ell_e=w_e R_{\rm eff}(e),
\]
where $R_{\rm eff}(e)$ is the effective resistance between the endpoints of the
edge in the weighted network.

The exact census in the accompanying script enumerates the two-terminal
series--parallel trees on six edges and solves for weights with $F=1/6$.
It finds the following six configurations.

\begin{center}
\begin{tabular}{p{0.38\linewidth}p{0.22\linewidth}p{0.18\linewidth}p{0.08\linewidth}}
\hline
Graph expression & weights & leverages & $|\cA|$\\
\hline
$P(S(e,e,e),e,e,e)$
& $1,1,1,\frac59,\frac59,\frac59$
& $(\frac{13}{18})^3,(\frac5{18})^3$ & 10\\
$P(S(e,e),S(e,e),e,e)$
& $1,1,1,1,\frac58,\frac58$
& $(\frac{11}{18})^4,(\frac5{18})^2$ & 12\\
$P(S(P(e,e,e),e),S(e,e))$
& $1,1,1,\frac95,\frac95,\frac95$
& $(\frac5{18})^3,(\frac{13}{18})^3$ & 10\\
$P(S(P(S(e,e),e,e),e),e)$
& $1,1,\frac58,\frac58,1,1$
& $(\frac{11}{18})^4,(\frac5{18})^2$ & 12\\
$P(S(P(e,e),P(e,e)),S(e,e))$
& $1,1,1,1,\frac85,\frac85$
& $(\frac7{18})^4,(\frac{13}{18})^2$ & 12\\
$P(S(P(e,e),e),S(P(e,e),e))$
& $1,1,\frac85,1,1,\frac85$
& $(\frac7{18})^4,(\frac{13}{18})^2$ & 12\\
\hline
\end{tabular}
\end{center}

\begin{proposition}\label{prop:six-tight}
For each of the six projectors above, $P^2=P$, $\tr P=3$, and
\[
  F(P)=\frac16.
\]
Moreover every active $3\times3$ block has simple least eigenvalue.
\end{proposition}

\begin{proof}
The statements are finite exact checks.  The projector identities are rational.
For each triple $T$, the characteristic polynomial of $P_{TT}$ is computed
exactly; the active triples have least eigenvalue $1/6$, while the inactive
triples in this census are singular with least eigenvalue $0$.  Simplicity is
checked by verifying that the second eigenvalue of each active block is strictly
larger than $1/6$.
\end{proof}

\section{Sharpness certificates}

Let $P_0\in\cE$, where
\[
  \cE=\{P\in\Gr(3,6):F(P)=1/6\}.
\]
Assume that the active least eigenvalues at $P_0$ are simple.  Choose a basis
$B_1,\ldots,B_9$ of $T_{P_0}\Gr(3,6)$.  For each active triple $T$, let $v_T$ be
a least-eigenvector of $(P_0)_{TT}$ and set
\[
  L_{Tj}=v_T^T(B_j)_{TT}v_T.
\]
Then for tangent coordinates $z\in\R^9$,
\[
  F'(P_0;z)=\max_{T\in\cA(P_0)} (Lz)_T.
\]

\begin{definition}
We call $P_0$ \emph{sharp} if
\[
  \kappa(P_0)=\min_{\|z\|=1}\max_{T\in\cA(P_0)}(Lz)_T>0.
\]
\end{definition}

\begin{lemma}[Sharpness criterion]\label{lem:sharp-criterion}
Suppose $\rank L=9$ and there is a vector $\lambda$ with all entries strictly
positive such that
\[
  L^T\lambda=0.
\]
Then $P_0$ is sharp.
\end{lemma}

\begin{proof}
If $\max_T(Lz)_T\le0$, then
\[
  0=\lambda^TLz=\sum_T\lambda_T(Lz)_T.
\]
This is a sum of nonpositive terms with strictly positive coefficients, hence
every $(Lz)_T=0$.  Since $\rank L=9$, this forces $z=0$.  Thus
$\max_T(Lz)_T>0$ for all $z\ne0$, and compactness of the unit sphere gives a
positive minimum.
\end{proof}

For some of the projectors the positive multiplier is not the most convenient
certificate.  The following equivalent rational LP certificate was used as the
load-bearing test.  Maximize
\[
  \sum_T(-Lz)_T
\]
over the rational polytope
\[
  Lz\le0,\qquad \|z\|_\infty\le1.
\]
The feasible point $z=0$ has objective $0$, and all summands are nonnegative on
the feasible set.  If the exact optimum is $0$ and $\rank L=9$, then every
feasible $z$ satisfies $Lz=0$, hence $z=0$.  Since $Lz\le0$ is a cone, this proves
that the critical cone is trivial and $\kappa(P_0)>0$.

\begin{theorem}[Sharpness of the six graphic-matroid extremals]\label{thm:six-sharp}
All six projectors in Proposition~\ref{prop:six-tight} are sharp points of
$F$ on $\Gr(3,6)$.
\end{theorem}

\begin{proof}
The proof is by exact rational computation.  A direct orthonormal tangent basis
contains square roots.  These are eliminated as follows.  Put
\[
  \widetilde B=B_{\rm red}^T,\qquad G=\widetilde B^TW\widetilde B,\qquad
  S=\widetilde B\,G^{-1}\widetilde B^TW.
\]
Then $S$ is a rational idempotent and
\[
  P=W^{1/2}SW^{-1/2}.
\]
Because $W^{1/2}$ is diagonal, this similarity restricts to every principal
block, so $P_{TT}$ and $S_{TT}$ have the same spectrum.  The active
eigenvectors, tangent functionals, and the LP certificate can therefore be
written over $\mathbb Q$ after positive row scalings of $L$.  For all six cases
the exact computation gives $\rank L=9$ and exact LP optimum $0$.
\end{proof}

\section{A seventh extremal outside the graphic census}

Numerical descent from random starts repeatedly finds an equality point whose
sorted leverage pattern is
\[
  \left(\frac5{14},\frac5{14},\frac5{14},
        \frac9{14},\frac9{14},\frac9{14}\right).
\]
This pattern is absent from the six graphic-matroid projectors above.  The
projector was reconstructed exactly over $\mathbb Q(\sqrt2,\sqrt5)$.  Its
off-diagonal entries have magnitudes among
\[
  \frac5{14},\qquad
  \frac{\sqrt5}{21},\qquad
  \frac{5\sqrt5}{42},\qquad
  \frac{5\sqrt2}{42},\qquad
  \frac{\sqrt{10}}{14}.
\]

\begin{theorem}[The out-of-family extremal]\label{thm:seventh}
There is an exact projector $P\in\Gr(3,6)$ with leverage pattern
$(5/14)^{\times3},(9/14)^{\times3}$ such that
\[
  F(P)=\frac16.
\]
It has thirteen active triples and seven singular triples.  Every active least
eigenvalue is simple, and $P$ is a sharp point of $F$ on $\Gr(3,6)$.
\end{theorem}

\begin{proof}
The entries are reconstructed from the numerical projector by PSLQ and then
checked exactly in $\mathbb Q(\sqrt2,\sqrt5)$.  The projector identities
$P^T=P$, $P^2=P$, $\tr P=3$, the leverage pattern, and the active/singular
triple counts are exact finite checks.  The sharpness certificate uses the same
critical-cone LP as Theorem~\ref{thm:six-sharp}.  In this case an additional
positive column scaling of the tangent basis rationalizes $L$; positive diagonal
coordinate changes preserve both $\rank L$ and the property
\[
  \{z:Lz\le0\}=\{0\}.
\]
The resulting rational matrix has rank $9$ and exact LP optimum $0$.
\end{proof}

\begin{remark}
An exhaustive search over the two-terminal series--parallel trees with the
weight parametrization used for the six graphic projectors does not produce the
leverage pattern $(5/14)^{\times3},(9/14)^{\times3}$.  Thus the seventh point is
not merely a relabelling of the displayed graphic census.
\end{remark}

\section{Isolation and conditional finiteness}

\begin{lemma}[Sharp points are isolated]\label{lem:sharp-isolated}
Let $P_0\in\cE$ be sharp and assume the active least eigenvalues at $P_0$ are
simple.  Then there is a neighborhood $U$ of $P_0$ in $\Gr(3,6)$ such that
\[
  U\cap\cE=\{P_0\}.
\]
\end{lemma}

\begin{proof}
Near $P_0$, inactive triples remain below $1/6$ by continuity.  Since the active
least eigenvalues are simple, the active functions $g_T$ are analytic.  Hence,
for $P=P_0+r\dot P$ in a local chart and $\|\dot P\|=1$,
\[
  F(P)=\frac16+rF'(P_0;\dot P)+O(r^2)
\]
uniformly in $\dot P$.  Sharpness gives $F'(P_0;\dot P)\ge\kappa>0$ on the unit
sphere, so $F(P)>1/6$ for sufficiently small $r>0$.
\end{proof}

\begin{corollary}
The seven projectors in Theorems~\ref{thm:six-sharp} and~\ref{thm:seventh} are
isolated points of $\cE$.
\end{corollary}

\begin{theorem}[Conditional finiteness]\label{thm:conditional-finite}
If every point of $\cE$ is sharp, then $\cE$ is finite.
\end{theorem}

\begin{proof}
The set $\cE=F^{-1}(1/6)$ is closed in the compact manifold $\Gr(3,6)$, hence
compact.  If every point is sharp, Lemma~\ref{lem:sharp-isolated} makes $\cE$ a
compact discrete space.  A compact discrete space is finite.
\end{proof}

\begin{remark}
The hypothesis of Theorem~\ref{thm:conditional-finite} is not proved.  It is a
global statement over $\cE$, and sampling cannot discharge it.  The theorem is
nevertheless useful because it identifies a concrete local certificate whose
failure would point to a positive-dimensional equality family.
\end{remark}

\section{Active-set orbit census and numerical KKT screens}

The twenty triples of $[6]$ are acted on by the row-permutation group $S_6$.
Burnside's lemma gives the following numbers of active-subset orbits by
cardinality:
{\scriptsize
\[
\begin{array}{c|rrrrrrrrrrrrrrrrrrrrr}
|\cA|&0&1&2&3&4&5&6&7&8&9&10&11&12&13&14&15&16&17&18&19&20\\
\hline
\#&1&1&3&7&21&43&94&161&249&312&352&312&249&161&94&43&21&7&3&1&1.
\end{array}
\]
}
Thus there are $2136$ active-set orbits total, $579$ nonempty orbits with
$|\cA|\le8$, and $1556$ orbits with $|\cA|\ge9$.  The seven certified equality
projectors occupy four distinct active-set orbits.

For a selected active set $\cA_0$, the numerical screen minimizes residuals for
the following conditions:
\[
  g_T(P)=t\quad(T\in\cA_0),\qquad
  g_T(P)\le t\quad(T\notin\cA_0),
\]
together with the stationarity condition
\[
  0=\Pi_{T_P\Gr(3,6)}\left(\sum_{T\in\cA_0}\alpha_T\,dg_T(P)\right),
  \qquad \alpha_T\ge0,\quad \sum_T\alpha_T=1.
\]
For counterexample searches the level condition is $t<1/6$; for equality
searches it is $t=1/6$.  This is a heuristic least-squares screen, not a proof.

\subsection*{Observed numerical results}

\begin{center}
\footnotesize
\begin{tabular}{llrrp{2.5in}}
\hline
mode & active sizes & orbit reps & candidates & note\\
\hline
counterexample & $1,\ldots,4$ & 32 & 0 & no tripwire\\
equality & $1,\ldots,4$ & 32 & 0 & no low-active point\\
counterexample & $5,6$ & 137 & 0 & no tripwire\\
equality & $5,6$ & 137 & 0 & no low-active point\\
counterexample & $7,8$ & 410 & 0 & no tripwire\\
equality & $7,8$ & 410 & 1 & selected subset of known 12-active point\\
counterexample & $9$ & 312 & 0 & one near-threshold non-candidate\\
equality & $9$ & 312 & 0 & no equality stratum\\
counterexample & $10,\ldots,13$ & 1074 & 0 & MLCore run\\
equality & $10,\ldots,13$ & 1074 & 2 & selected subsets of known 12- and 13-active points\\
counterexample & $14,\ldots,20$ & 170 & 0 & no tripwire\\
equality & $14,\ldots,20$ & 170 & 0 & no equality stratum\\
\hline
\end{tabular}
\end{center}

The equality hit with selected $|\cA_0|=8$ had residual about $7.1\cdot10^{-13}$,
but its actual active set had size $12$ and matched a known TTSP active-set
orbit.  Thus it is not an eight-active equality stratum.

The two equality hits from selected sizes $10,\ldots,13$ were likewise
selected-subset hits.  Reconstructing the projectors from the numerical
parameters and recomputing the actual active sets gave one known TTSP
12-active orbit and the out-of-family 13-active orbit.  In both cases the
actual active-gradient matrix has rank $9$ and a strictly positive numerical
multiplier, so neither is a non-sharp candidate.

The best size-nine counterexample near-miss polished to
\[
  F=0.16666994562575838>1/6
\]
with actual top active size one.  The best MLCore record for selected
$|\cA_0|=12$ had
\[
  F=0.166677230>1/6
\]
and again actual top active size one.  No numerical run found a point with
$F<1/6-10^{-6}$.

Across all equality-mode screens, the actual-active-set classifier found three
low-residual equality hits, all in known sharp active-set orbits; it found no new
actual active-set orbit, no actual active size at most $9$, and no numerical
non-sharp signature.

\begin{remark}
These screens are useful mainly as triage.  They also expose a bookkeeping
hazard: a selected subset of a larger actual active set may satisfy equality and
stationarity to high precision.  Any exact active-set proof must therefore track
the actual active set, not only the selected equations.
\end{remark}

\section{A semialgebraic finiteness target}

The conditional finiteness theorem reduces finiteness of $\cE$ to excluding
non-sharp equality points.  In the simple-active case this obstruction has a
small polynomial representation.  Work in the Grassmann chart
\[
  Y=\begin{pmatrix}I_3\\ Z\end{pmatrix},\qquad
  P(Z)=Y(Y^TY)^{-1}Y^T=\frac{N(Z)}{d(Z)},\qquad d=\det(Y^TY)>0.
\]
For a triple $T$, put
\[
  M_T(Z)=6N(Z)_{TT}-d(Z)I_3.
\]
Then $\lambda_{\min}(P(Z)_{TT})=1/6$ with simple least eigenvalue is encoded by
\[
  M_T(Z)\succeq0,\qquad \det M_T(Z)=0,\qquad \rank M_T(Z)=2.
\]
The raw determinant has the universal factor $d^2$; since this chart has
$d>0$, the active determinant equation may equivalently be written as the
saturated degree-six equation
\[
  \frac{\det M_T(Z)}{d(Z)^2}=0.
\]
In purely algebraic computations this open chart condition should be retained,
for instance by adjoining an inverse variable $u$ and the equation
$u\,d(Z)-1=0$.
On a cofactor patch, a nonzero kernel vector $w_T(Z)$ is obtained as the cross
product of two rows of $M_T$.  If $h\in\R^9$ is a nonzero tangent witness for
non-sharpness, then after clearing the positive denominator $d^2$ one obtains the
polynomial inequalities
\[
  w_T^T\left[
  \sum_j h_j\bigl(d\,\partial_jN-N\,\partial_jd\bigr)_{TT}
  \right]w_T\le0,\qquad T\in\cA,
\]
with the normalization $\sum_jh_j^2=1$.

Thus the simple-active non-sharp obstruction lives in only 18 variables: nine
chart variables and nine tangent variables.  The active-set size changes the
number of equations and inequalities, but not the ambient dimension.  A degree
audit gives $\deg d=6$, entries of $N$ of degree at most $6$, raw active
determinant degree at most $18$ but saturated active determinant degree $6$,
cofactor-patch norm degree at most $24$, and cleared directional inequalities
of degree at most $35$.

There is also a useful count.  At a simple-active equality point, sharpness means
that the active gradients positively span the 9-dimensional tangent space.  This
requires at least ten active gradients.  Hence a simple-active equality point
with $|\cA|\le9$ would automatically be non-sharp.  The equality screens found no
such point, and all observed equality points have $|\cA|\in\{10,12,13\}$.

The nonsimple case, where an active least eigenvalue has multiplicity greater
than one, must be split off separately: one active triple then contributes a
set-valued subgradient rather than one row of the active-gradient matrix.

\section{A CAD cascade for the base active set}\label{sec:cad-cascade}

The semialgebraic target becomes more concrete on the active set of the base
graphic extremal $P(S(e,e,e),e,e,e)$.  We use zero-based row labels and write
\[
  \cA_0=\{012,013,014,015,023,024,025,123,124,125\},
\]
where, for instance, $013$ denotes the triple $\{0,1,3\}$.  Let
$I_0\subset\mathbb Q[z_0,\ldots,z_8,u]$ be generated by the ten saturated active
determinants
\[
  \frac{\det(6N(Z)_{TT}-d(Z)I_3)}{d(Z)^2},\qquad T\in\cA_0,
\]
together with the chart equation $u\,d(Z)-1=0$.  For an active triple $T$ and
two local row indices $E\subset\{0,1,2\}$, let $\Delta_{T,E}$ denote the
corresponding principal $2\times2$ minor of $6N(Z)_{TT}-d(Z)I_3$.

The following three polynomial relations were found and then checked by
open-locus Groebner computations:
\begin{align*}
  R_0&=\Delta_{013,\{0,2\}}+46656z_0^2,\\
  R_1&=\Delta_{013,\{1,2\}}+46656z_1^2,\\
  R_2&=\Delta_{014,\{0,2\}}+46656z_7^2.
\end{align*}
On a real point satisfying active positive semidefiniteness, each displayed
minor is nonnegative.  Therefore $R_i=0$ forces the corresponding square
coordinate to vanish.

\begin{proposition}[Computed cascade on the base stratum]\label{prop:cad-cascade}
Over $\mathbb Q$ in the full-rank chart $d>0$, the following statements hold.
\begin{enumerate}[leftmargin=*]
\item The determinant ideal $I_0$ has dimension $3$.  The open locus
$I_0+\langle aR_0-1\rangle$ is zero-dimensional of degree $608$.  Hence $R_0$
vanishes on every positive-dimensional irreducible component of $V(I_0)$.
\item Put $I_1=I_0+\langle R_0,z_0\rangle$.  Then
$\dim I_1=2$ and the top degree is $16$.  Moreover
$I_1+\langle aR_1-1\rangle$ is empty, so $R_1$ holds on this whole branch.
\item Put $I_2=I_1+\langle R_1,z_1\rangle$.  Then
$\dim I_2=1$ and the top degree is $24$.  The normal form of $R_2$ modulo
$I_2$ is zero, so $R_2\in I_2$.
\item Put $I_3=I_2+\langle R_2,z_7\rangle$.  Then $I_3$ is zero-dimensional of
degree $56$, and a lexicographic basis is
\[
\begin{gathered}
z_2^2-5,\quad 216u-1,\quad z_8^2,\quad z_7,\quad z_6^2-5,\quad z_5^2,\\
z_4^2-5,\quad z_3z_5z_8,\quad z_3^2,\quad z_1,\quad z_0.
\end{gathered}
\]
Its real radical is therefore
\[
z_0=z_1=z_3=z_5=z_7=z_8=0,\qquad
z_2^2=z_4^2=z_6^2=5,\qquad u=1/216.
\]
The resulting eight real sign choices all fail the GTZ semialgebraic
conditions: an active block has $\lambda_{\min}-1/6=-1/6$, and an inactive block
has least eigenvalue $2/3$ above the threshold.
\item The finite open residue from the first step is also accounted for.  Let
$J_0=I_0+\langle aR_0-1\rangle$.  The computed $z_0$ eliminant of $J_0$ is
\[
\begin{split}
&z_0^{10}-\frac{815}{441}z_0^8-\frac{93850}{64827}z_0^6
  +\frac{2295875}{453789}z_0^4\\
&\hspace{25mm}
  -\frac{13883125}{4084101}z_0^2+\frac{8556250}{12252303}\\
&=(z_0^2-\tfrac59)(z_0^2-\tfrac{10}{21})(z_0^2+\tfrac53)
  \left(z_0^4-\frac{365}{147}z_0^2+\frac{34225}{21609}\right).
\end{split}
\]
The last two factors have no real zeros.  On the factor $z_0^2=5/9$, an exact
lex basis reduces to
\[
\begin{gathered}
9z_1^2-5,\quad 25980a-1,\quad 6u-1,\quad
9z_8^2-5,\quad 9z_7^2-5,\quad 9z_6^2-5,\quad 9z_5^2-5,\\
5z_4-9z_5z_7z_8,\quad
5z_3-9z_5z_6z_8,\quad
5z_2-9z_7z_8z_1,\quad
5z_0-9z_6z_7z_1 .
\end{gathered}
\]
This gives $32$ real sign choices; direct substitution shows that all have the
same active set $\cA_0$ and satisfy the sharp equalities and inequalities.
On the factor $z_0^2=10/21$, the $z_1$ eliminant factors as
\[
(z_1^2-\tfrac{10}{21})(z_1^2+\tfrac53)
\left(z_1^4-\frac{365}{147}z_1^2+\frac{34225}{21609}\right).
\]
The only real subfactor is $z_1^2=10/21$, and adjoining it gives a lex basis
whose first polynomial is $3z_2^2+5$.  Hence this branch has no real points.
\end{enumerate}
\end{proposition}

\begin{proof}
The algebraic dimension and degree statements are Groebner-basis computations
using Singular over $\mathbb Q$ with its modular backend.  The open-locus tests
adjoin a fresh inverse variable $a$ and the equation $aR_i-1=0$; if this ideal
has smaller dimension, or is empty, then $R_i$ cannot be nonzero on the
corresponding component.  The positive-semidefinite forcing step is elementary:
on the active semialgebraic locus the minors $\Delta_{T,E}$ are nonnegative, and
so are the displayed squares.  The forced finite lex basis gives the eight real
points explicitly, and direct exact substitution checks the active and inactive
eigenvalue inequalities.  The open-residue eliminant and the displayed
real-factor refinements were computed by exact characteristic-zero Groebner
computations; a separate six-prime CRT reconstruction gives the same $z_0$
eliminant.
\end{proof}

\begin{remark}
Proposition~\ref{prop:cad-cascade} is not yet a proof of the full
$(6,3)$ case.  It accounts for the positive-dimensional residual determinant
branches and the finite $R_0\ne0$ residue attached to the base active set in the
standard chart.  The size-twelve and thirteen-active chart computations below
propagate the same finiteness program to the other known equality active-set
orbits in this standard chart.  Other active-set orbits and chart patches remain.
\end{remark}

\section{A finite chart for the thirteen-active point}\label{sec:seventh-chart}

The out-of-family extremal of Theorem~\ref{thm:seventh} gives another exact
chart computation.  Let
\[
  \cA_7=\{012,014,015,023,024,025,034,035,123,124,125,134,135\}.
\]
Let $I_7$ be the saturated determinant ideal for these active triples, again in
the standard chart $Y=[I;Z]$ with $u\,d(Z)-1=0$.

\begin{proposition}[Computed finite chart for the seventh point]
The ideal $I_7$ is zero-dimensional of degree $32$ over $\mathbb Q$.  A
lexicographic basis is
\[
\begin{gathered}
8z_0^2-5,\quad 63u-8,\quad 8z_8^2-5,\quad
64z_7^2-45,\quad 64z_6^2-45,\quad 8z_5^2-5,\\
5z_4-8z_5z_7z_8,\quad
5z_3-8z_5z_6z_8,\quad z_2,\quad
45z_1-64z_6z_7z_0 .
\end{gathered}
\]
Consequently the real chart points form a $32$-point sign family.  Direct
substitution shows that all $32$ points satisfy the GTZ inequalities with
$F=1/6$, and their actual active set is exactly $\cA_7$.
\end{proposition}

\begin{proof}
The displayed lex basis is the output of an exact Singular computation over
$\mathbb Q$ using the modular backend.  It leaves five independent signs, hence
$32$ real points.  The active and inactive inequalities are then checked by
substituting the displayed formulas into the $3\times3$ principal blocks.
\end{proof}

\section{Finite real branches for the two size-twelve charts}
\label{sec:size-twelve-charts}

Two known equality active-set orbits have twelve active triples in the standard
chart:
\[
\begin{aligned}
  \cA_{12a}
    &=\{012,013,023,024,025,034,035,123,124,125,134,135\},\\
  \cA_{12b}
    &=\{012,013,014,015,024,025,034,035,124,125,134,135\}.
\end{aligned}
\]
Let $I_{12a}$ and $I_{12b}$ be the corresponding saturated determinant ideals
with $u\,d(Z)-1=0$.

\begin{proposition}[Computed finite real branches for the size-twelve strata]
For each of $I_{12a}$ and $I_{12b}$, an exact Groebner basis contains
\[
  729u^3+81u^2-9u-1=729\left(u-\frac19\right)\left(u+\frac19\right)^2 .
\]
Since $d(Z)=\det(Y^TY)>0$ on the real full-rank chart and $u=1/d(Z)$, every
real chart point lies on the branch $u=1/9$.  On this branch both ideals are
zero-dimensional of degree $32$ over $\mathbb Q$.

For $\cA_{12a}$, a lexicographic basis of the real branch is
\[
\begin{gathered}
z_0^2-1,\quad 9u-1,\quad z_8,\quad 8z_7^2-5,\quad 8z_6^2-5,\quad z_5,\quad
8z_4^2-5,\\
5z_3-8z_4z_6z_7,\quad z_2^2-1,\quad 5z_1-8z_6z_7z_0 .
\end{gathered}
\]
For $\cA_{12b}$, a lexicographic basis of the real branch is
\[
\begin{gathered}
z_2^2-1,\quad 9u-1,\quad z_8^2-1,\quad 8z_7^2-5,\quad 8z_6^2-5,\quad
z_5^2-1,\\
z_4-z_5z_7z_8,\quad z_3-z_5z_6z_8,\quad z_1,\quad z_0 .
\end{gathered}
\]
Thus each real branch consists of a $32$-point sign family.  Direct substitution
checks that every point in each family satisfies the GTZ inequalities with
$F=1/6$, and that the actual active set is exactly the selected twelve-triple
set.
\end{proposition}

\begin{proof}
The branch polynomial in $u$ is read from exact characteristic-zero Groebner
bases.  The displayed lex bases are exact computations after adding $9u-1$.
The sign counts and active/inactive inequalities are verified by substituting
the resulting formulas into the projector $P(Z)=Y(Y^TY)^{-1}Y^T$.
\end{proof}

\section{A structured low-active branch}
\label{sec:structured-low-active}

The active-set census also leaves low-active determinant strata that are much
harder than the known equality strata.  The most persistent size-six survivor in
our searches has
\[
  \cA_6=\{012,013,045,145,234,235\}.
\]
A refined over-tied numerical root has five outside triples tied at a common
level:
\[
  \{024,034,125,135,245\}.
\]
In the standard chart the root has the pattern
\[
Z=
\begin{pmatrix}
-a&a&-b\\
c&d&-a\\
-d&-c&-a
\end{pmatrix}.
\]
Substituting this ansatz into the six active determinant equations and five
outside-tie determinant equations gives a zero-dimensional system of degree
\(280\), both over \(\mathbb F_{32003}\) and over \(\mathbb Q\).  The exact
\(\mathbb Q\) grevlex basis contains the simple tail relations
\[
  a^2((c+d)^2-5)=0,\qquad
  ((c+d)^2-5)((c-d)^2-5)=0.
\]
Consequently every real point of the ansatz system lies either on \(a=0\) or
on \((c+d)^2=5\).  The \(a=0\) branch is zero-dimensional of degree \(96\);
an exact FGLM conversion gives a six-polynomial lex basis, and Sage
enumeration over the algebraic real field finds no real points.  It remains to
handle \((c+d)^2=5\), which is the branch containing the refined numerical
root.

On this branch, an exact \(\mathbb Q\) grevlex computation gives quotient
degree \(184\), matching the modular computation.  Modular lex computations over
six primes have stable quotient degree \(184\) and a degree-\(22\) eliminant in
the tie parameter \(q\), where the tied outside triples have least eigenvalue
\(q/6\).  The visible factor
\[
  (q-3)(q-5)^2(q^2-3q+6)^2(q^2-9q+24)^2
\]
removes nuisance branches.  The six-prime CRT reconstruction of the residual
factor is
\[
\begin{split}
&q^{11}-\frac{575}{16}q^{10}+\frac{562849}{1024}q^9
-\frac{4768527}{1024}q^8+\frac{97733649}{4096}q^7\\
&\quad-\frac{311529979}{4096}q^6+\frac{154321449}{1024}q^5
-\frac{374696325}{2048}q^4+\frac{543960225}{4096}q^3\\
&\quad-\frac{225237375}{4096}q^2+\frac{24046875}{2048}q
-\frac{253125}{256},
\end{split}
\]
which factors over \(\mathbb Q\) as a cubic times an octic:
\[
\begin{split}
&\left(q^3-9q^2+\frac{81}{4}q-\frac{45}{4}\right)\\
&\quad\cdot
\left(q^8-\frac{431}{16}q^7+\frac{293857}{1024}q^6
-\frac{776859}{512}q^5+\frac{2094513}{512}q^4
-\frac{1353145}{256}q^3\right.\\
&\hspace{27mm}\left.
+\frac{3301125}{1024}q^2-\frac{453375}{512}q
+\frac{5625}{64}\right).
\end{split}
\]
Multiplying this residual factor by the visible nuisance factor and reducing
against the exact \(\mathbb Q\) grevlex basis of the branch gives zero normal
form.  Thus the degree-\(22\) \(q\)-polynomial is certified in characteristic
zero; the CRT step is no longer only modular evidence.

Converting the exact branch basis to lexicographic order by FGLM gives quotient
dimension \(184\), again matching the grevlex and modular computations.  After
adjoining the residual factors separately, the cubic branch has lex quotient
dimension \(24\) and no real points: its lex basis contains \(b^2+1\) and
\(a^2\).  The octic branch has lex quotient dimension \(32\) and contains
\(16\) real points.  The nuisance factors \(q=3\) and \(q=5\) have lex quotient
dimension \(32\) each, with \(8\) and \(16\) real points respectively.

The refined numerical root corresponds to the real octic root
\[
  q=1.180695836766579\ldots,
\]
so the tied inactive triples sit above the threshold by
\[
  \frac{q-1}{6}=0.030115972794\ldots.
\]

\begin{theorem}[Exclusion of the structured size-six ansatz]
There is no real point in the four-parameter structured ansatz above which
realizes the selected threshold pattern: the six selected triples at threshold
with \(P_{TT}-I/6\succeq0\), and every inactive triple satisfying
\(\lambda_{\min}(P_{TT})\le 1/6\).
\end{theorem}

\begin{proof}
The tail relation \(a^2((c+d)^2-5)=0\) splits the ansatz system into the
\(a=0\) branch and the \((c+d)^2=5\) branch.  The \(a=0\) branch has an exact
zero-dimensional lex basis of quotient dimension \(96\), and exact enumeration
over the algebraic real field gives no real points.

On the \((c+d)^2=5\) branch, the exact \(\mathbb Q\) branch ideal has quotient
dimension \(184\), and its \(q\)-eliminant is the certified degree-\(22\)
polynomial above.  The non-real quadratic nuisance factors have no real points.
The cubic residual factor has no real points because the exact lex basis on
that component contains \(b^2+1\).  It remains to inspect the octic residual
factor and the real nuisance factors \(q=3\) and \(q=5\).

For these three lex systems, Sage enumeration over the algebraic real field
finds \(16\), \(8\), and \(16\) real points respectively.  At each point we
compute the exact numerator
\[
  N=Y\,\operatorname{adj}(Y^TY)\,Y^T,\qquad d=\det(Y^TY)>0,
\]
so that \(P_{TT}-I/6\succeq0\) is equivalent to
\[
  6N_{TT}-dI\succeq0 .
\]
For an active triple, a negative principal minor of this shifted block excludes
the point.  For an inactive triple, a positive-definite shifted block certifies
\(\lambda_{\min}(P_{TT})>1/6\), also excluding a point from the desired
threshold pattern.

The exact sign certificate classifies all \(40\) real points.  On the octic
component, \(12\) points have a negative active principal minor and \(4\) points
have an inactive positive-definite shifted block.  On the \(q=3\) component,
all \(8\) real points fail active positive semidefiniteness.  On the \(q=5\)
component, all \(16\) real points have an inactive positive-definite shifted
block.  There are no unclassified real points.
\end{proof}

A deterministic numerical classifier using the lifted \(q\)-values on the
\((c+d)^2=5\) branch found the same \(40\)-point real-root profile in two
independent runs.  This is now only a diagnostic cross-check; the exclusion
above is the exact real-algebraic certificate for the structured ansatz.

\begin{remark}
This theorem excludes only the displayed four-parameter structured ansatz for
the selected size-six active and outside-tie pattern.  It does not prove the
full inequality \eqref{eq:gtz63}.
\end{remark}

\section{What remains}

Proposition~\ref{prop:kkt} isolates the remaining difficulty: prove that no
Clarke-stationary point of $F$ has value below $1/6$.  The exact local analysis
above, including the CAD cascade in Proposition~\ref{prop:cad-cascade} and the
finite branch computations in Sections~\ref{sec:seventh-chart}
and~\ref{sec:size-twelve-charts}, and the structured low-active evidence in
Section~\ref{sec:structured-low-active}, does not solve that global problem.

The most realistic continuation is a semialgebraic exclusion of the non-sharp
obstruction just described.  For each active-set stratum, modulo row-permutation
symmetry and cofactor-patch choices, one can ask an exact real-algebraic solver
to certify that the equality, positivity, and tangent-witness conditions have no
solution outside the known sharp active-set orbits.  Inactive constraints may be
relaxed for an infeasibility proof: if the relaxed obstruction is empty, then the
true obstruction is empty.

A second, more direct but currently less effective, continuation is certified
global optimization over $\Gr(3,6)$:
\begin{enumerate}[leftmargin=*]
\item Use the exact sharpness certificates to remove explicit neighborhoods of
the seven known equality points.  A second-order eigenvalue bound gives a radius
of the form
\[
  F(P_0+r\dot P)\ge \frac16+\kappa r-\frac{r^2}{\operatorname{gap}_{\min}},
\]
where $\operatorname{gap}_{\min}$ is the smallest active spectral gap.
\item Cover the remaining compact set by interval boxes in a Grassmann chart or
by several charts, using interval bounds for the twenty $3\times3$ least
eigenvalues.
\item Branch on active triples near nonsmooth regions.  The blocks are only
$3\times3$, so interval characteristic-polynomial or Rayleigh quotient bounds
are feasible.
\end{enumerate}
This route is computational rather than elegant, but it is honest: it targets the
global $\forall$ statement directly and does not rely on extrapolating from
sampled extremals.

A relaxed bound $F(P)\ge 1/6-\varepsilon$ is also a useful intermediate target,
but only as part of a gap-closing strategy.  By itself it would not imply the
sharp hypothesis.  What it can do is separate the problem into a certified
compact region with a margin gap and a small near-threshold region that must be
handled by the exact equality and KKT machinery above.

\section{Reproducibility}\label{sec:repro}

The accompanying repository contains the following verification scripts in
\texttt{verify/}.
\begin{center}
\small
\begin{tabular}{@{}ll@{}}
\texttt{v4\_nesterenko.py} & exact tightness of a graphic extremal\\
\texttt{v6\_sharpness.py} & exact sharpness certificate for the base extremal\\
\texttt{v9\_rational\_certificate.py} & six graphic-matroid sharpness certificates\\
\texttt{v11\_seventh\_exact.py} & exact reconstruction and sharpness of the seventh point\\
\texttt{v13\_radius.py} & computed local radii at the seven points\\
\texttt{v16\_jacobian.py} & independent Jacobian-rank isolation check\\
\texttt{v17\_active\_orbits.py} & active-set orbit counts\\
\texttt{v18\_kkt\_screen.py} & numerical active-set KKT screens\\
\texttt{v20\_finiteness\_probe.py} & actual-active-set classifier for equality hits\\
\texttt{v21\_semialgebraic\_route.py} & semialgebraic obstruction sizing\\
\texttt{v35\_relaxed\_one\_ninth.py} & exact constants in the \(F(P)\ge1/9\) proof
\end{tabular}
\end{center}
The CAD cascade in Section~\ref{sec:cad-cascade} is supported by the following
Sage/Singular helpers in \texttt{code/sage/}.
\begin{center}
\small
\begin{tabular}{@{}ll@{}}
\texttt{probe\_cad\_obstruction\_strata.py} & forced-coordinate dimension drops\\
\texttt{lexify\_determinant\_slice.py} & exact slices and lex bases\\
\texttt{classify\_lex\_slice\_roots.py} & real-root and inequality classifier\\
\texttt{screen\_lex\_slice\_minors.py} & minor-relation discovery on slices\\
\texttt{probe\_cad\_relation\_open\_locus.py} & component-relation checks\\
\texttt{probe\_cad\_relation\_normal\_forms.py} & ideal-membership checks\\
\texttt{reconstruct\_univariate\_crt.py} & modular eliminant reconstruction\\
\texttt{dump\_determinant\_gb.py} & exact determinant Groebner bases\\
\texttt{dump\_facstd\_components.py} & modular component diagnostics\\
\texttt{probe\_s6\_ansatz.py} & structured size-six ansatz probes\\
\texttt{verify\_s6\_ansatz\_eliminant.py} & exact \(q\)-eliminant normal-form check\\
\texttt{lexify\_s6\_ansatz\_basis.py} & exact FGLM conversion of ansatz branches\\
\texttt{certify\_s6\_ansatz\_signs.py} & exact real sign checks on ansatz branches\\
\texttt{classify\_s6\_ansatz\_numeric.py} & numerical classification of ansatz roots
\end{tabular}
\end{center}
The principal output artifacts are:
\begin{quote}
\small\raggedright
\path{code/sage/out/cad_open_R0_on_I_QQ.json}\\
\path{code/sage/out/cad_open_R1_on_z0_QQ.json}\\
\path{code/sage/out/cad_open_R2_on_z0z1_QQ.json}\\
\path{code/sage/out/cad_obstruction_known_base_invd_QQ_z0z1z7_relz1_relz7.json}\\
\path{code/sage/out/lex_cad_R0_nonzero_QQmod_z0.json}\\
\path{code/sage/out/lex_cad_R0_nonzero_z0_eliminant_crt_6primes.json}\\
\path{code/sage/out/lex_cad_R0_nonzero_QQmod_z0sq_5_9_z1.json}\\
\path{code/sage/out/lex_cad_R0_nonzero_QQmod_z0sq_10_21_z1.json}\\
\path{code/sage/out/lex_cad_R0_nonzero_QQmod_z0z1sq_10_21_z2.json}\\
\path{code/sage/out/lex_det_mask32253_QQmod_z0.json}\\
\path{code/sage/out/gb_det_mask32243_QQ.json}\\
\path{code/sage/out/gb_det_mask31215_QQmod.json}\\
\path{code/sage/out/lex_det_mask32243_QQmod_u_pos_z0.json}\\
\path{code/sage/out/lex_det_mask31215_QQmod_u_pos_z2.json}\\
\path{code/sage/out/s6_ansatz_azero_QQ.json}\\
\path{code/sage/out/s6_ansatz_azero_fglm_QQ.json}\\
\path{code/sage/out/s6_ansatz_cplusd_QQ.json}\\
\path{code/sage/out/s6_ansatz_cplusd_q_residual_crt_6primes.json}\\
\path{code/sage/out/s6_ansatz_cplusd_eliminant_verify_QQ.json}\\
\path{code/sage/out/s6_ansatz_cplusd_fglm_QQ.json}\\
\path{code/sage/out/s6_ansatz_cplusd_octic_fglm_QQ.json}\\
\path{code/sage/out/s6_ansatz_cplusd_q3_fglm_QQ.json}\\
\path{code/sage/out/s6_ansatz_cplusd_q5_fglm_QQ.json}\\
\path{code/sage/out/s6_ansatz_cplusd_sign_cert_QQ.json}\\
\path{code/sage/out/s6_ansatz_structured_sign_cert_QQ.json}\\
\path{code/sage/out/classify_s6_ansatz_cplusd_numeric_s800.json}
\end{quote}
The exact claims in this paper are reproducible by running
\[
\begin{split}
&\texttt{python verify/v4\_nesterenko.py},\\
&\texttt{python verify/v9\_rational\_certificate.py},\\
&\texttt{python verify/v11\_seventh\_exact.py},\\
&\texttt{python verify/v17\_active\_orbits.py}.
\end{split}
\]
The numerical KKT sweep over active sizes $10,\ldots,13$ was run on MLCore.
\begin{center}
\small
\begin{tabular}{@{}lll@{}}
job & project & region\\
\texttt{gtz63-kkt-x79fea} & \texttt{test-paos} & \texttt{ix-m5-sm12}
\end{tabular}
\end{center}

\begin{thebibliography}{9}

\bibitem{GTZ}
S.~A. Goreinov, E.~E. Tyrtyshnikov, and N.~L. Zamarashkin.
\newblock A theory of pseudoskeleton approximations.
\newblock \emph{Linear Algebra and its Applications}, 261:1--21, 1997.

\bibitem{GTZ2}
S.~A. Goreinov, E.~E. Tyrtyshnikov, and N.~L. Zamarashkin.
\newblock Pseudo-skeleton approximations by matrices of maximal volume.
\newblock \emph{Mathematical Notes}, 62:515--519, 1997.

\bibitem{SenguptaPautov}
R.~Sengupta and M.~Pautov.
\newblock Submatrices with the best-bounded inverses: revisiting the hypothesis.
\newblock arXiv:2604.05944, 2026.

\bibitem{DerezinskiWarmuth}
M.~Derezi\'nski and M.~K. Warmuth.
\newblock Unbiased estimates for linear regression via volume sampling.
\newblock arXiv:1705.06908, 2018.

\bibitem{HornJohnson}
R.~A. Horn and C.~R. Johnson.
\newblock \emph{Matrix Analysis}.
\newblock Cambridge University Press, second edition, 2013.

\end{thebibliography}

\end{document}
